% =====================================================================
%  PART I — AMBUJA CEMENTS LTD
% =====================================================================
\partpage{Part I}{Ambuja Cements Ltd}{The difficult-to-value name. A company that
has bought a great deal of cement capacity, earns very little on it, and is priced
as though it already earns a great deal.}

\runheadlogo{logo_amb}{Part I \textbar\ Ambuja Cements}
\setcounter{section}{0}

\banner{Ambuja Cements Ltd}{NSE: AMBUJACEM \textbar\ BSE: 500425}{SELL}{Rs 99.4}{Rs 400.90}{$-$75.2\%}{Consolidated}

\begin{paracol}{2}
\section{Investment case and recommendation}

\lead{We initiate on Ambuja Cements with a SELL and an intrinsic value of Rs 99.4
per share against a traded price of Rs 400.90.} Five discounted cash flow methods
— FCFF, capital cash flow, adjusted present value, FCFE and economic value added —
reconcile to that figure to the last paisa. The gap to the market is not a rounding
question or a discount-rate question. It is a question about returns.

The central finding is arithmetic rather than rhetorical. Return on invested
capital including goodwill is 4.6\% in FY2026-27 and reaches only 8.7\% by
FY2035-36, against a weighted average cost of capital of 11.36\%. Economic value
added is negative in nine of the ten forecast years. When a company earns less on
capital than capital costs, enterprise value sits below invested capital, and that
is precisely what this model produces.

\subsection{Three things explain the return}

\subsubsection{The plants are underworked} Capacity utilisation is 67.9\% against
82\% at UltraTech. Ambuja owns 109 million tonnes of capacity and sold 74 million
tonnes. A large asset base is being carried by a smaller volume than it was built
for, and fixed costs are spread thinly as a result.

\subsubsection{Unit economics are weak} EBITDA per tonne is Rs 886. UltraTech's
core assets earn Rs 966 and the wider peer set earns more. The gap is concentrated
inside the standalone parent, which runs an 11.7\% EBITDA margin while the
subsidiaries run 23.3\%.

\subsubsection{A third of the capital was bought, not built} Acquisitions left
Rs 22,979 crore of goodwill and other intangibles on the balance sheet — 34.2\% of
invested capital. That is capital shareholders paid for and must earn a return on.
Strip goodwill out and ROIC improves from 4.6\% to 7.0\%, which measures the plants
but flatters what shareholders actually earned.

\begin{keybox}\boxtitle{The one number that matters}
To justify Rs 400.90 the market must underwrite a terminal return on invested
capital comfortably above 15\%, sustained in perpetuity. Ambuja has not earned
above its cost of capital in either year of consolidated history the annual report
provides.
\end{keybox}

\subsection{What we already give the company}

This is not a bearish forecast dressed as analysis. The model assumes utilisation
climbs from 67.9\% to 85\%, EBITDA per tonne rises 70\% over ten years from Rs 886
to Rs 1,502, the standalone margin more than doubles from 11.7\% to 19.0\% as the
Sanghi and Penna integrations complete, and capacity grows from 109 to 133 million
tonnes. On those assumptions the value is Rs 99.4. The forecast is generous and the
answer is still a quarter of the price.

\switchcolumn
\ex{amb_football.png}{\textbf{Exhibit A1.} Ranges are the model's own outputs; peer
multiples are applied to Ambuja's consolidated FY2025-26 EBITDA, normalised PAT and
net worth. Source: authors' model; peer annual reports; NSE close 30 March 2026.}

\ex{amb_roic.png}{\textbf{Exhibit A4.} NOPAT is EBIT taxed at the statutory 25.168\%
throughout; the reported effective rate is negative and is never used. Invested
capital is measured on opening balances. Source: authors' model.}

\ex{amb_eva.png}{\textbf{Exhibit A5.} EVA is NOPAT less a capital charge of WACC on
opening operating invested capital. It turns positive only in the final forecast
year. Source: authors' model.}

\switchcolumn*
\section{Company profile: what it sells, where it operates, who owns it}

\subsection{The company}

Ambuja Cements Ltd was incorporated in 1983 as Gujarat Ambuja Cements and became
one of India's largest cement producers over the following three decades. Control
passed to Holcim in stages from 2006, and in September 2022 the Adani group acquired
Holcim's Indian cement assets — Ambuja together with its listed subsidiary ACC — in
what remains the largest acquisition in Indian building materials. The registered
office moved to Adani Corporate House, Ahmedabad.

Since that change of control the group has expanded by acquisition rather than
principally by greenfield construction: Sanghi Industries and Penna Cement were
absorbed into the standalone entity by amalgamation, and Orient Cement was acquired
from the CK Birla group. Consolidated capacity at the valuation date was 109 MTPA
against a stated target of 119 MTPA, and consolidated cement sales volume was
74 million tonnes on revenue of Rs 40,656 crore.

\begin{keybox}\boxtitle{Why the acquisition history is the valuation}
Four acquisitions in three years is why Rs 22,979 crore of goodwill and intangibles
sits on the balance sheet, why a third of invested capital was bought rather than
built, and why utilisation is 67.9\%. The corporate history and the return on capital
are the same story told twice.
\end{keybox}

\subsection{Products and brands}

The group is a single-product business in economic terms: grey cement is the
overwhelming majority of revenue. What varies is the brand and the blend. Ordinary
Portland Cement, Portland Pozzolana Cement and Portland Slag Cement are sold under
the Ambuja, ACC and Orient names, with premium and application-specific variants —
Ambuja Plus, Ambuja Kawach for water resistance, Ambuja Compocem, and ACC's Gold and
Suraksha ranges — carrying higher realisation per tonne. Ready-mix concrete and
aggregates, largely through ACC Concrete, make up the small remainder.

For valuation purposes the mix matters less than the tonne. The model is built on
capacity times utilisation times realisation per tonne precisely because that is what
the business actually is, and because premium-blend share is one of the few levers
management holds over realisation in a market where price has been flat to falling.

\subsection{Board of directors}

The board is chaired by \textbf{Gautam Adani} in a non-executive, non-independent
capacity, with \textbf{Karan Adani} also non-executive and non-independent.
\textbf{Vinod Bahety} is Wholetime Director and Chief Executive Officer, and is the
executive accountable for the integration programme on which this valuation turns.
Five independent directors — \textbf{Rajnish Kumar}, former Chairman of the State Bank
of India, \textbf{Maheswar Sahu}, \textbf{Purvi Sheth}, \textbf{Ameet Desai} and
\textbf{Praveen Garg} — make up 62.5\% of the board. Photographs, roles and the
governance analysis follow on the next page.

\begin{okbox}\boxtitleg{Sourcing note}
The board list, the committee structure, the 62.5\% independence ratio and the 96.24\%
attendance figure are taken from the Board of Directors pages of the FY2025-26
Integrated Annual Report (pp. 230--231). The shareholding pattern is from the
Shareholder Information section of the same report and the company's
shareholding-pattern filings. Nothing in this section comes from a financial data
aggregator.
\end{okbox}

\subsection{Operating footprint}

The group operates a national network of integrated plants, grinding units and
terminals. Consolidated capacity was 109 MTPA at the valuation date against a stated
target of 119 MTPA, with 10.2 MTPA of scheduled expansion due to complete by
FY2026-27. Rs 9,085 crore of capital work in progress was still to be commissioned —
which is why free cash flow is negative in the first two forecast years and why
utilisation, not capacity, is the binding constraint.

Management's own stated priorities for the period map onto the same three levers this
report measures: complete the scheduled expansion; integrate the acquired assets
across logistics, procurement and market execution; and improve the energy mix,
circular materials and lower-clinker product share while preserving cost
competitiveness.

\subsection{Shareholders and shareholding pattern}

Ambuja is a promoter-controlled company with a substantial free float. The Adani
promoter group held \textbf{67.64\%} of the equity share capital at 31 March 2026 —
67.68\% on a voting-rights basis — leaving a public float of 32.36\%. On 10 April 2026
the company issued 1,29,93,708 equity shares to eligible shareholders of Sanghi
Industries under the amalgamation, which diluted the promoter holding to 67.29\%
(67.33\% on voting rights).

Two features of the register matter for the valuation. First, the free float is large
enough that the traded price is a genuine market clearing price rather than a thin
quote, so the 75\% gap this report identifies cannot be dismissed as an artefact of
illiquidity. Second, the pending merger of ACC, Orient, Sanghi and Penna into Ambuja
raises the share count by roughly 13.7\% and further dilutes the promoter stake — the
mechanism by which minority leakage is removed, and the reason both a pre-merger and a
post-merger per-share figure are presented in the valuation section.

\switchcolumn
\ex{amb_ownership.png}{\textbf{Exhibit A17.} Promoter percentage is company-disclosed
at 31 March 2026. ACC remained separately listed at the valuation date; all four
entities are being merged into Ambuja. Source: company disclosures; FY2025-26 annual
report, subsidiary listing.}

\ex{amb_products.png}{\textbf{Exhibit A18.} Grey cement is the business; the brand and
the blend decide realisation per tonne. The company does not publish a full revenue
split by product, so the split shown is indicative of the disclosed mix. Source:
FY2025-26 annual report; authors' analysis.}

\ex{amb_ownership_donut.png}{\textbf{Exhibit A20.} Promoter shares are reported
unpledged. Domestic institutions own 17.5\% between mutual funds and other
institutions; foreign institutional ownership is modest at 5.6\%, so the register is
domestically anchored. Source: FY2025-26 Integrated Annual Report, Shareholder
Information; shareholding-pattern filings.}

\ex{amb_shareholding.png}{\textbf{Exhibit A19.} The promoter/float split at the
valuation date itself. Source: company disclosure of the Sanghi amalgamation share
issue.}

\ex{amb_plantmap.png}{\textbf{Exhibit A25.} National plant, grinding-unit and terminal
footprint. Cement is freight-constrained, so regional density matters more than
national share. Source: FY2025-26 Integrated Annual Report, At a Glance.}

\begin{extab}{The company at a glance}{Source: FY2025-26 Integrated Annual Report;
NSE close 30 March 2026.}
\begin{tabular}{@{}p{0.44\linewidth}r@{}}
\hrow \wh{Item} & \wh{Value}\\
Incorporated & 1983\\
\arow Promoter group & Adani\\
Control acquired & September 2022\\
\arow Registered office & Ahmedabad\\
Listings & NSE, BSE\\
\arow Capacity, consolidated (MTPA) & 109\\
Capacity target, stated (MTPA) & 119\\
\arow Sales volume FY26 (Mn.T) & 74.0\\
Capacity utilisation FY26 & 67.9\% \\
\arow Consolidated revenue (Rs cr) & 40,656\\
Consolidated EBITDA (Rs cr) & 6,559\\
\arow EBITDA per tonne (Rs) & 886\\
Shares outstanding (crore) & 247.18\\
\arow Market capitalisation (Rs cr) & 99,095\\
Promoter holding, 31-Mar-26 & 67.64\% \\
\arow Public float & 32.36\% \\
Listed subsidiary & ACC Ltd (50.05\%)\\
\end{tabular}
\end{extab}

\begin{extab}{Board metrics}{Source: FY2025-26 Integrated Annual Report,
Board of Directors, pp. 230--231.}
\begin{tabular}{@{}p{0.62\linewidth}r@{}}
\hrow \wh{Metric} & \wh{Value}\\
Directors on the board & 8\\
\arow of which independent & 5\\
\trow \textbf{Independent directors} & \textbf{62.5\%}\\
Average attendance, board and committees & 96.24\% \\
\arow Age profile, 35--55 years & 37.5\% \\
Age profile, 56--75 years & 62.5\% \\
\arow Statutory committees & 5\\
Additional governance committees & 6\\
\end{tabular}
\end{extab}

\end{paracol}
\clearpage

% ---------------- full-width board and management spread ----------------
\fullhead{Board of directors}{Ambuja Cements Ltd, as constituted at 31 March 2026.
Source: FY2025-26 Integrated Annual Report, Board of Directors, pp. 230--231.}

\noindent
\person{0.113\textwidth}{amb_gautam_adani.png}{Gautam Adani}{Non-executive Chairman, Non-independent}\hfill
\person{0.113\textwidth}{amb_karan_adani.png}{Karan Adani}{Non-executive, Non-independent Director}\hfill
\person{0.113\textwidth}{amb_vinod_bahety.png}{Vinod Bahety}{Wholetime Director \& CEO}\hfill
\person{0.113\textwidth}{amb_rajnish_kumar.png}{Rajnish Kumar}{Non-executive, Independent Director}\hfill
\person{0.113\textwidth}{amb_maheswar_sahu.png}{Maheswar Sahu}{Non-executive, Independent Director}\hfill
\person{0.113\textwidth}{amb_purvi_sheth.png}{Purvi Sheth}{Non-executive, Independent Director}\hfill
\person{0.113\textwidth}{amb_ameet_desai.png}{Ameet Desai}{Non-executive, Independent Director}\hfill
\person{0.113\textwidth}{amb_praveen_garg.png}{Praveen Garg}{Non-executive, Independent Director}

\vspace{6mm}
\begin{paracol}{2}
\subsection*{Composition and coverage}

Five of the eight directors are independent, so \textbf{independent directors
comprise 62.5\% of the board}. Reported average attendance at board and committee
meetings was \textbf{96.24\%} in FY2025-26. The age profile splits 37.5\% in the
35--55 band and 62.5\% in the 56--75 band.

Statutory committees cover audit, nomination and remuneration, stakeholder
relationships, corporate social responsibility and risk management. Beyond the
statutory minimum the board runs committees for mergers and acquisitions, technology
and data security, commodity prices, legal and tax, and corporate responsibility and
reputation risk. That is unusually broad committee coverage for an Indian
mid-to-large-cap industrial.

\begin{keybox}\boxtitle{Formal governance is strong; the question is capital allocation}
Breadth of committee coverage is not the issue at Ambuja. The substantive governance
question this report raises is narrower and harder: whether the board is testing
acquisition synergies, utilisation ramps and impairment assumptions against
\emph{realised cash returns}. Four acquisitions in three years produced Rs 22,979
crore of goodwill and a return on capital that has not reached its cost. The critical
bridge is from transaction approval to post-acquisition accountability — capex,
integration cost, minority dilution and ROIC tracked together, not separately.
\end{keybox}

\switchcolumn
\subsection*{Where independent challenge should bite}

Independent directors chair or sit on the audit, risk, nomination, technology, legal
and responsibility committees. On the evidence in this report, their most material
contribution is ensuring that reported scale, premiumisation and ESG progress are
reconciled with free cash flow and with the return on the \emph{full} capital base —
including the goodwill that EBITDA cannot see.

\begin{warnbox}\boxtitler{Three questions a shareholder should put to the board}
\textbf{1.} What return on invested capital, including acquisition goodwill, does the
board hold management to, and over what period?\\[2pt]
\textbf{2.} At what point does a utilisation ramp that has not materialised trigger an
impairment review of the Rs 22,979 crore carried as goodwill and intangibles?\\[2pt]
\textbf{3.} On the merger of ACC, Orient, Sanghi and Penna, what is the board's own
estimate of value created net of the 13.7\% dilution — expressed in cash, not in EPS?
\end{warnbox}
\end{paracol}

\vspace{3mm}
\fullhead{Senior management}{The leadership team executing the integration on which
the bull case depends. Source: FY2025-26 Integrated Annual Report, Leadership Team,
pp. 232--233.}

\noindent
\person{0.150\textwidth}{amb_l_vinod_bahety.png}{Vinod Bahety}{Wholetime Director \& CEO}\hfill
\person{0.150\textwidth}{amb_rohit_soni.png}{Rohit Soni}{Chief Financial Officer}\hfill
\person{0.150\textwidth}{amb_sanjay_gupta.png}{Sanjay Kumar Gupta}{Chief Procurement Officer}\hfill
\person{0.150\textwidth}{amb_praveen_k_garg.png}{Praveen Kumar Garg}{Chief Logistics Officer}\hfill
\person{0.150\textwidth}{amb_bhimsi_kachhot.png}{Bhimsi Kachhot}{Chief Strategy \& Business Development}

\vspace{3mm}
\noindent
\person{0.150\textwidth}{amb_sanjay_behl.png}{Sanjay Behl}{Head -- Sales, Marketing \& Logistics}\hfill
\person{0.150\textwidth}{amb_john_varghese.png}{John Varghese}{Chief People Officer}\hfill
\person{0.150\textwidth}{amb_madhavi_isanaka.png}{Madhavi Isanaka}{Chief Digital Officer}\hfill
\person{0.150\textwidth}{amb_vineet_bose.png}{Vineet Bose}{Chief Legal Officer}\hfill
\person{0.150\textwidth}{amb_vaibhav_dixit.png}{Vaibhav Dixit}{Head -- Manufacturing}

\vspace{4mm}
\begin{paracol}{2}
\textbf{Execution priorities.} The scorecard that matters for this valuation is
narrow: fill the plants, lift EBITDA per tonne toward the peer set, convert the
Sanghi and Penna assets from a drag into contributors, and complete the merger
without leaking value through the exchange ratios. Four of the ten roles above —
procurement, logistics, manufacturing and strategy — map directly onto the cost per
tonne that decides whether return on capital ever crosses the cost of capital.

\switchcolumn
\textbf{What to watch.} Consolidated EBITDA per tonne is the single disclosure that
separates this valuation from the market price. Below it, the two operating levers
are utilisation (67.9\% today against 82\% at UltraTech) and the blended cost per
tonne that procurement, logistics and manufacturing jointly control. Premium mix
helps realisation at the margin but cannot by itself close a 75\% valuation gap.
\end{paracol}
\clearpage

\begin{paracol}{2}
\section{Margin stack, working capital and depreciation}

Three FY2025-26 disclosures explain why a better cost performance did not reach the
operating line. Taken together they are the thesis of this report compressed into one
financial year.

\subsection{The direct economics improved}

Manufacturing gross margin — revenue less materials, stock-in-trade, inventory
movement and power and fuel — rose from 57.9\% to 58.8\%. Delivered gross margin,
which also deducts freight, improved from 34.4\% to 35.4\%, with freight broadly flat
at 23.4\% of revenue. On the direct cost line, FY2025-26 was a good year.

\subsection{But EBITDA margin fell}

EBITDA margin nevertheless declined from 16.9\% to 16.1\%, because employee and other
expenses rose from 17.5\% to 19.3\% of revenue. The saving made on materials and fuel
was consumed by overhead, much of it integration-related.

\subsection{And EBIT margin fell much further}

EBIT margin dropped from 10.4\% to 7.4\%. That 3.1-point fall is larger than the
EBITDA decline, and the difference is depreciation. D\&A rose 55.4\% to Rs 3,570
crore, from 6.5\% to 8.8\% of revenue and from 38.4\% to 54.4\% of EBITDA. Within it,
intangible amortisation more than doubled to Rs 607 crore as acquired brands, mining
rights and other intangibles entered the consolidated base.

\begin{keybox}\boxtitle{This is the thesis in one financial year}
Ambuja got better at making cement in FY2025-26 and worse at earning a return. The
gain was real and it was captured at the gross line. It was then absorbed by
integration overhead and, decisively, by a depreciation charge rising nearly four
times faster than revenue. A margin-recovery case must show where employee and other
expenses normalise; it cannot simply assume price improvement.
\end{keybox}

\subsection{Working capital: real, but borrowed}

Core trade-cycle working capital fell from Rs 2,843 crore to Rs 2,236 crore, releasing
about Rs 608 crore, and the cash conversion cycle improved from 32.7 to 20.8 days.
Read the composition, though. Inventory days fell 4.1 and receivable days were flat,
while payable days rose 7.6 as payables grew 39.7\%. The improvement came mostly from
\textbf{supplier funding}, not from faster collection. Within inventory the mix also
shifted: fuel inventory fell 26.8\% while stores and spares rose 46.8\%,
work-in-progress 51.3\% and finished goods 30.4\% — consistent with commissioning and
integration activity rather than with tighter management.

\begin{okbox}\boxtitleg{Why we did not capitalise the improvement}
Continued reliance on higher payable days is not an operating margin improvement and
should not be extrapolated. Our forecast holds net working capital at minus 0.34\% of
revenue — essentially the FY2025-26 actual — so no cash is manufactured from the
working-capital line in any forecast year. Note also that receivables ageing
deteriorated beneath a stable headline: the 6--12 month bucket rose to Rs 111 crore
from Rs 30 crore.
\end{okbox}

\switchcolumn
\ex{amb_marginstack.png}{\textbf{Exhibit A21.} Manufacturing gross margin deducts
materials, stock-in-trade, inventory movement and power and fuel; delivered gross
margin also deducts freight. Neither is a reported Ind AS subtotal, so both are
computed from the disclosed expense notes. Source: FY2025-26 consolidated financial
statements, Notes 44--51.}

\ex{amb_da.png}{\textbf{Exhibit A22.} The composition of the charge. Intangible
amortisation is the fastest-growing element and sits on a balance the forecast holds
flat. Source: FY2025-26 consolidated financial statements, Note 49.}

\ex{amb_ccc.png}{\textbf{Exhibit A23.} DIO and DPO use cash operating cost; DSO uses
revenue from operations. All day metrics are closing-balance proxies, because a
comparable opening FY2024-25 group balance is not available after restatement and
acquisitions. Source: FY2025-26 consolidated financial statements, Notes 15, 17 and
35.}

\begin{warnbox}\boxtitler{This resolves an audit item, against us}
Our model charges depreciation at 8.0\% of opening \emph{operating} fixed assets while
holding the Rs 9,433 crore intangible balance flat and never amortising it. Historical
D\&A includes Rs 607 crore of amortisation on exactly those intangibles. The rate and
the asset base do not cover the same assets — the asset-consistency PARTIAL recorded
in Appendix C. We can now put a number on it. Modelling the amortisation explicitly
would reduce forecast NOPAT and \emph{lower} the valuation further, so the omission
runs against our own recommendation.
\end{warnbox}

\end{paracol}
\clearpage

% ---------------- full-width product spread ----------------
\fullhead{Products and brands}{Core and premium cement, and allied building
materials. Product images from the FY2025-26 Integrated Annual Report and the Ambuja
product portfolio.}

\noindent
\product{0.185\textwidth}{amb_core.png}{Adani Ambuja Cement}{The core high-performance product for general residential and commercial construction.}\hfill
\product{0.185\textwidth}{amb_plus.png}{Ambuja Plus}{Formulated for stronger, denser concrete. Sits in the premium range.}\hfill
\product{0.185\textwidth}{amb_kawach.png}{Ambuja Kawach}{Water-repellent formulation, positioned as a super-premium product.}\hfill
\product{0.185\textwidth}{amb_compocem.png}{Ambuja Compocem}{Slag and silicate enriched, for a finer, brighter, lower-impact finish.}\hfill
\product{0.185\textwidth}{amb_coolwalls.png}{Ambuja Cool Walls}{Eco-friendly walling blocks designed to keep homes up to 5 degrees cooler.}

\vspace{6mm}
\begin{paracol}{2}
\subsection*{Why the portfolio matters, and how much}

Cement is a commodity and differentiation is limited to brand, blend and service.
Ambuja nonetheless runs a real premium ladder: Plus for strength, Kawach for moisture
resistance, Compocem for finish and lower clinker intensity. Premium and specialised
products accounted for roughly \textbf{35\% of trade sales} in FY2025-26. Four
products — Ambuja Cement, Plus, Kawach and Compocem — carry GRIHA certification as
low-impact products, which matters for institutional and government demand.

Alongside the bagged range sit Cool Walls, an eco-friendly walling block, and
Alccofine, a specialised supplementary cementitious material for demanding
applications. ACC Concrete supplies ready-mix.

\switchcolumn
\subsection*{The commercial reading}

Mix improvement supports realisation and brand differentiation, and it is a genuine
lever — but it is a second-order one. On our forecast, realisation grows below
inflation for most of the ten-year period and roughly 60\% of the assumed improvement
in EBITDA per tonne comes from cost and fixed-cost absorption rather than from price.

\begin{okbox}\boxtitleg{The honest weighting}
A premium ladder that reaches 35\% of trade sales is a good business asset. It is not
large enough to change the answer. The consolidated EBITDA outcome still depends on
plant utilisation and the cost base, which is why the valuation is built on capacity
times utilisation times EBITDA per tonne rather than on product mix.
\end{okbox}
\end{paracol}
\clearpage
\begin{paracol}{2}
\section{How the group is put together for valuation}

Ambuja Cements is the Indian cement arm of the Adani group. At the valuation date
the group comprised the standalone parent — which itself absorbed Sanghi Industries
and Penna Cement by amalgamation — together with a 50.05\% holding in the separately
listed ACC Ltd and majority holdings in Orient Cement and others. Consolidated
capacity was 109 MTPA against a stated target of 119 MTPA, and consolidated cement
sales volume was 74 million tonnes on revenue of Rs 40,656 crore.

\subsection{Why this valuation is consolidated}

Ambuja standalone carries 61.6\% of group revenue but only 44.7\% of group EBITDA.
The parent runs an 11.7\% margin; the subsidiaries run 23.3\%. The weak assets —
Sanghi and Penna — sit inside the standalone entity, while the better assets, ACC
and Orient, sit outside it. A standalone model would value the worse half of the
business and ignore the better half, while the market capitalisation prices the
whole group.

\begin{warnbox}\boxtitler{A category error worth avoiding}
Dividing the group market capitalisation of Rs 99,095 crore by \emph{standalone}
EBITDA of Rs 2,930 crore prints an EV/EBITDA of 34.0x and is meaningless. Adding
the non-controlling interest of Rs 12,499 crore to enterprise value and dividing by
\emph{consolidated} EBITDA of Rs 6,559 crore gives 17.0x. That is the pairing used
throughout this report.
\end{warnbox}

\subsection{The forecast is a build-up, not a single margin}

Consolidated revenue is built from physical capacity times utilisation times
realisation per tonne, which disciplines the forecast against the plant the company
actually owns. That consolidated revenue is then split into the standalone parent
and the subsidiary block using the FY2025-26 revenue share of 61.6\%, and each block
carries its own margin path. Consolidated EBITDA is the sum of the two blocks, never
an assumed group margin.

\subsection{The merger changes the share count, not the economics}

ACC, Orient, Sanghi and Penna are being merged into Ambuja. Pre-merger, the
non-controlling interest of Rs 12,499 crore is deducted in the bridge and the value
is divided by today's 247.18 crore shares. Post-merger, the minority is acquired and
the share count rises roughly 13.7\%. Both routes should give a similar answer, and
they do: Rs 99.4 pre-merger against Rs 132 post-merger, the difference being the
discount at which minorities are being bought in.

\begin{notebox}\boxtitlea{Open item}
The swap ratios and the 13.7\% dilution reached this model through broker research,
which our sourcing policy does not admit. They must be re-sourced from the scheme
filing before the post-merger figure is relied upon.
\end{notebox}

\switchcolumn
\ex{amb_entity.png}{\textbf{Exhibit A15.} The subsidiary column is derived as
consolidated less standalone and therefore ties to the reported total by
construction. Intercompany eliminations sit inside the subsidiary column because
they cannot be separated from the published statements. Source: FY2025-26 annual
report, consolidated and standalone statements of profit and loss.}

\ex{amb_margins.png}{\textbf{Exhibit A8.} The parent starts at 11.7\% because Sanghi
and Penna sit inside it; the subsidiaries start at 23.3\% because ACC and Orient sit
outside it. The consolidated line is the weighted outcome, not an input. Source:
authors' model.}

\ex{amb_capacity.png}{\textbf{Exhibit A6.} Volume is built as capacity times
utilisation rather than assumed as a growth rate, so projected volume can never
exceed the plant the company owns. Source: authors' model; capacity target of 119
MTPA per the FY2025-26 annual report.}

\switchcolumn*
\section{Industry context, demand and decarbonisation}

Indian cement is a regional, freight-constrained commodity. Cement cannot travel far
economically, so competition is decided within a few hundred kilometres of a plant and
national market share means less than regional density. Demand is driven by housing,
infrastructure and government capital expenditure.

Three features of the current cycle matter for this valuation. First, the industry has
been adding capacity faster than demand, which holds utilisation near 80\% and caps
pricing power. Second, realisation has been flat to falling in real terms:
UltraTech's realisation compounded at 2.0\% over ten years and fell 6.1\% in
FY2024-25. Third, consolidation has been aggressive, and the acquirers have paid
prices that put goodwill on balance sheets — which is exactly the problem this report
measures.

\begin{keybox}\boxtitle{Why the sector's own multiple cannot settle the argument}
A peer median EV/EBITDA of 19.4x, solved back through the steady-state identity,
implies perpetual growth of 10.99\%. Nominal GDP on our own assumptions is 10.76\%.
A company cannot outgrow the economy forever, so the peer multiple is the
implausible end of the range, not the benchmark against which our DCF should be
corrected.
\end{keybox}

\subsection{The global backdrop}

The IMF projects global growth of 3.0\% in 2026 and 3.4\% in 2027. The economy is
resilient but uneven: technology investment supports activity while geopolitical and
energy shocks weigh on importers, and disinflation has stalled — leaving interest-rate
and energy-cost volatility directly relevant to cement demand and fuel procurement.

Indian cement is predominantly a domestic industry, but the transmission is
asymmetric: imported coal, petcoke, freight and financing conditions transmit global
shocks into margins faster than global construction volumes transmit into demand. That
asymmetry is why a cost-led margin thesis is more fragile than a volume-led one.

\subsection{The Indian market}

India is the world's second-largest cement producer. Industry production reached
approximately \textbf{491.4 million tonnes} in FY2025-26, up 8.6\%, and ICRA's
FY2026-27 volume outlook is 7--8\%. Infrastructure, urban housing and industrial capex
remain the structural demand engines. Installed capacity is projected to reach roughly
850 MTPA by 2030, a 5.1\% compound rate from 2025.

\begin{keybox}\boxtitle{Demand is not the problem}
Volume growth of 7--8\% against nominal GDP of 10.76\% is a perfectly healthy
industry. The difficulty is that capacity is being added at least as fast, which holds
utilisation near 80\% across the sector and keeps pricing competitive. Earnings
improvement therefore has to come from utilisation, integration and fixed-cost
absorption rather than from sustained real price increases — which is exactly what
both forecasts in this report assume.
\end{keybox}

\subsection{Decarbonisation, and why it is a cost story}

Global cement production has softened in recent years but direct emissions intensity
has not materially improved. Decarbonisation depends on lower clinker factors,
supplementary cementitious materials, alternative fuels, material efficiency and
eventually carbon capture. The IEA estimates that early near-zero cement plants using
carbon capture can cost 75--150\% more than conventional production, so green-premium
demand and policy support remain important to the economics.

For both companies in this report, blended products and renewable power lower
emissions \emph{and} operating cost simultaneously. Carbon capture is a longer-dated
capital and policy question that does not bear on a valuation struck today.

\begin{warnbox}\boxtitler{Why ESG does not enter this valuation separately}
We do not adjust the discount rate or the cash flows for ESG factors, and this is a
deliberate methodological position rather than an omission. The green energy and waste
heat recovery programmes already appear in both forecasts as the cost improvement that
lifts EBITDA per tonne. Adding an ESG premium or discount on top would count the same
effect twice — a straightforward breach of the assumptions consistency this report is
built on.

We also do not use third-party ESG scores. They vary widely between providers for the
same company, their methodologies are not disclosed in enough detail to audit, and
they are not an authorised source under the data policy applied throughout this work.
\end{warnbox}

\switchcolumn
\ex{ind_india.png}{\textbf{Exhibit A24.} Industry production and the two companies'
directional share of it. Sources: IBEF, Indian Cement Industry, May 2026; ICRA volume
outlook; company reports.}

\begin{extab}{India cement industry indicators}{Sources: IBEF, May 2026; ICRA;
company reports.}
\begin{tabular}{@{}p{0.60\linewidth}r@{}}
\hrow \wh{Indicator} & \wh{Value}\\
FY2025-26 production & 491.4 Mn.T\\
\arow FY2025-26 growth & 8.6\% \\
FY2026-27 volume outlook & 7--8\% \\
\arow 2030 installed capacity projection & 850 MTPA\\
2025--30 capacity CAGR & 5.1\% \\
\arow Sector utilisation, approximate & c.80\% \\
\trow Ambuja share of output & \textbf{c.15\%}\\
\trow UltraTech share of output & \textbf{c.29\%}\\
\end{tabular}
\end{extab}

\begin{okbox}\boxtitleg{The one number that ties industry to valuation}
If the industry adds capacity at 5.1\% a year while demand grows 7--8\%, utilisation
improves only slowly and pricing stays competitive. That is the environment both
models assume: realisation growing at or below inflation, with margin improvement
coming from cost and absorption. A reader who expects a genuine pricing cycle should
say so explicitly and re-run the sensitivity grid — the tables are built to take it.
\end{okbox}

\begin{extab}{Global context}{Sources: IMF, July 2026 World Economic Outlook Update;
IEA, Breakthrough Agenda Report 2025 — Cement and Concrete.}
\begin{tabular}{@{}p{0.62\linewidth}r@{}}
\hrow \wh{Indicator} & \wh{Value}\\
Global growth, 2026 & 3.0\% \\
\arow Global growth, 2027 & 3.4\% \\
Near-zero cement cost premium (IEA) & 75--150\% \\
\arow Ambuja green power ambition & rising\\
UltraTech green power share & 43\% \\
\arow UltraTech green power ambition & 65\% \\
\end{tabular}
\end{extab}

\switchcolumn*
\section{Historical analysis and value drivers}

\subsection{Two years of consolidated history, and why}

The FY2025-26 annual report provides consolidated statements for FY2025-26 and a
restated FY2024-25, and no more. Earlier consolidated years exist only in broker
research, which is not an authorised source. A return trend asserted from two points
is thin, and we say so rather than manufacture depth. Two further distortions matter:
FY2024-25 was restated for the Sanghi and Penna amalgamation, and the transition year
to March 2023 covered fifteen months. No growth rate should be taken off the reported
series without adjustment.

\subsection{The tax line is not usable as reported}

Ambuja's reported consolidated tax line for FY2025-26 is a net \emph{credit} of
Rs 2,338 crore, so reported profit exceeds profit before tax and the reported
effective rate is minus 71\%. Every return measure in this report is therefore
computed at the statutory rate of 25.168\% under section 115BAA. The credit
decomposes into three drivers, and treating each on its economics rather than its
accounting turns a headline gain into a loss of Rs 11.54 per share.

\switchcolumn
\ex{amb_tax.png}{\textbf{Exhibit A13.} Driver 1 is excluded because the model taxes
EBIT at the statutory rate throughout, so the reversal was never in the cash flows.
Driver 1a is real cash already paid and recoverable, and sits outside operating
invested capital. Driver 2 is deducted once rather than amortised. Driver 3 is memo
only because it is already netted inside Driver 2. Source: FY2025-26 annual report,
notes 32 and 53; authors' analysis.}

\begin{extab}{The three drivers of the Rs 2,338 crore credit}{Source: FY2025-26
annual report, note 32(ii) and note 53.}
\begin{tabular}{@{}p{0.40\linewidth}r p{0.30\linewidth}@{}}
\hrow \wh{Driver} & \wh{Rs cr} & \wh{Treatment}\\
\arow 1. Litigation provisions released & 1,251 & Excluded\\
1a. Tax paid under protest & 614 & Added as asset\\
\arow 2. Deferred tax liabilities & 3,466 & Deducted once\\
3. s.72A loss shelter & 346 & Memo only\\
\trow \textbf{Net effect on value} & & \textbf{$-$Rs 11.54/sh}\\
\end{tabular}
\end{extab}

\begin{extab}{Historical returns, consolidated}{NOPAT is EBIT at the statutory
25.168\%. Source: FY2025-26 annual report; authors' analysis.}
\begin{tabular}{@{}p{0.46\linewidth}rr@{}}
\hrow \wh{} & \wh{FY25R} & \wh{FY26}\\
Revenue incl. grants (Rs cr) & 35,336 & 40,656\\
\arow EBITDA (Rs cr) & 5,984 & 6,559\\
EBITDA margin & 16.9\% & 16.1\% \\
\arow EBITDA per tonne (Rs) & n/a & 886\\
Capacity utilisation & n/a & 67.9\% \\
\arow ROIC incl. goodwill & 5.2\% & 3.3\% \\
ROIC excl. goodwill & 7.6\% & 5.0\% \\
\arow Goodwill \% of invested capital & 30.9\% & 34.2\% \\
Reported effective tax rate & 13.3\% & $-$70.9\% \\
\end{tabular}
\end{extab}

\switchcolumn*
\subsection{The value drivers, ranked}

Three levers move this valuation, and they are the three the forecast narrative
spends its words on.

\subsubsection{Capacity utilisation} The single most important driver. The plants
exist and are paid for; whether they fill decides whether the fixed cost base is
absorbed. Every point of utilisation is close to pure margin.

\subsubsection{EBITDA per tonne} The metric the industry actually steers by. Our
path takes it from Rs 886 to Rs 1,502, with roughly 60\% of the improvement assumed
to come from cost and fixed-cost absorption rather than price.

\subsubsection{The capital base} Goodwill of Rs 22,979 crore does not shrink and
does not earn a separate return. It is the denominator problem that no operating
improvement of plausible size can outrun.

\begin{okbox}\boxtitleg{What the two ROIC lines are for}
Excluding goodwill, the business earns a mediocre return that measures the plants.
Including goodwill, the return measures what shareholders actually earned on what
they actually paid. Both are shown throughout. The second is the one that decides
the valuation.
\end{okbox}

\section{Forecast assumptions}

The explicit period is ten years, not the customary five. Ambuja runs at 67.9\%
utilisation and earns far below its cost of capital today; a five-year window ends
before the ramp has played out and throws almost the entire value into the terminal
period. A longer explicit horizon is the standard treatment for a difficult-to-value
company, and it holds terminal value to 72\% of enterprise value rather than the
85\%-plus a five-year horizon would produce.

Every driver is set out in the table opposite with the anchor it came from. Two
deserve comment. Realisation is assumed to grow \emph{below} inflation for most of
the forecast, consistent with the sector's record and with sell-side assumptions of
flat to falling net sales realisation. And net working capital is held at minus
0.34\% of revenue, essentially the FY2025-26 actual, so no cash is manufactured by
shrinking working capital.

\begin{warnbox}\boxtitler{The assumption that carries the case}
The standalone EBITDA margin rising from 11.7\% to 19.0\% is the single most
important assumption in the model. It embeds a complete turnaround of the Sanghi and
Penna assets. If it does not happen, the value is materially lower than Rs 99.4.
\end{warnbox}

\switchcolumn
\begin{extab}{Forecast drivers, FY2026-27 to FY2035-36}{Every figure is an input on
the Assumptions sheet of the model. Source: authors' assumptions anchored to the
FY2025-26 annual report.}
\begin{tabular}{@{}p{0.44\linewidth}rrr@{}}
\hrow \wh{Driver} & \wh{FY27} & \wh{FY31} & \wh{FY36}\\
Capacity (MTPA) & 112 & 123 & 133\\
\arow Utilisation & 70.0\% & 79.5\% & 85.0\% \\
Volume (Mn.T) & 78.4 & 97.8 & 113.1\\
\arow Realisation growth & 4.8\% & 2.0\% & 1.2\% \\
Standalone revenue share & 61.5\% & 60.7\% & 60.0\% \\
\arow Standalone EBITDA margin & 13.0\% & 16.7\% & 19.0\% \\
Subsidiary EBITDA margin & 23.6\% & 25.0\% & 26.0\% \\
\arow Consolidated margin (output) & 17.1\% & 20.0\% & 21.8\% \\
EBITDA per tonne (Rs) & 983 & 1,288 & 1,502\\
\arow Capex (Rs cr) & 9,500 & 7,500 & 7,500\\
Depreciation, \% opening assets & 8.0\% & 8.0\% & 8.0\% \\
\arow NWC, \% of revenue & $-$0.34\% & $-$0.34\% & $-$0.34\% \\
\end{tabular}
\end{extab}

\ex{amb_ebitdat.png}{\textbf{Exhibit A7.} The endpoint of Rs 1,502 sits 56\% above
what UltraTech's best assets currently earn. The assumption is deliberately
demanding: we wanted the valuation to fail on the capital base rather than on a
conservative operating forecast. Source: authors' model.}

\ex{amb_ic.png}{\textbf{Exhibit A9.} Goodwill is held flat in nominal terms across
the forecast. It stays inside invested capital for the return measure, because it is
capital shareholders paid, but it is excluded from terminal reinvestment because
organic growth does not require buying more of it. Source: authors' model.}

\switchcolumn*
\section{Cost of capital}

Ambuja carries almost no debt — borrowings and leases of Rs 866 crore against a
market capitalisation of Rs 99,095 crore, a debt-to-equity ratio of 0.87\%. The
capital structure machinery therefore barely binds: WACC, the unlevered rate $\rho$
and the cost of equity sit within four basis points of each other.

\subsection{Inputs}

The risk-free rate is the FBIL 10-year G-Sec par yield of 6.85\% for the week ended
20 March 2026, the last weekly print before the valuation date, less the India
sovereign default spread of 1.87\% — because the equity risk premium already
contains a country risk premium and double-counting it would be a risk-consistency
error. That gives a default-free rupee rate of 4.98\%. The equity risk premium of
7.08\% and the unlevered building-materials beta of 0.90 are Damodaran, January 2026;
the cash-corrected beta is used because cash is added back separately at the bridge.

\subsection{Relevering: the debt policy has to match the formula}

The model assumes a constant debt-to-value ratio with continuous rebalancing —
Harris--Pringle, Case 2 in the course taxonomy — which is what allows the interest
tax shield to be discounted at $\rho$ and what makes all five methods agree. Under
that policy the relevering formula carries \emph{no} tax term: $\beta_L =
\beta_U(1+D/E)$. This report applies that formula consistently, giving a levered
beta of 0.9079 and a cost of equity of 11.41\%.

\begin{okbox}\boxtitleg{Circularity}
At a debt-to-equity ratio of 0.87\%, no plausible revision to the equity value moves
the weights enough to change WACC in the second decimal. The circularity problem is
immaterial here, and we say so having checked it rather than having assumed it.
\end{okbox}

\switchcolumn
\begin{extab}{Cost of capital, Ambuja}{Harris--Pringle constant-rebalancing policy
applied consistently in the unlever, the relever and the shield. Source: RBI/FBIL;
Damodaran January 2026; authors' analysis.}
\begin{tabular}{@{}p{0.60\linewidth}r@{}}
\hrow \wh{Input} & \wh{Value}\\
10-year G-Sec par yield (FBIL, w/e 20-Mar-26) & 6.85\% \\
\arow Less: India sovereign default spread & $-$1.87\% \\
\trow \textbf{Risk-free rate (default-free rupee)} & \textbf{4.98\%} \\
India equity risk premium & 7.08\% \\
\arow Unlevered beta, building materials & 0.900\\
Market debt / equity & 0.87\% \\
\arow Levered beta, $\beta_U(1+D/E)$ & 0.9079\\
\trow \textbf{Cost of equity (CAPM)} & \textbf{11.41\%} \\
Pre-tax cost of debt & 8.00\% \\
\arow Statutory tax rate (s.115BAA) & 25.168\% \\
Debt / enterprise value & 0.87\% \\
\trow \textbf{WACC} & \textbf{11.36\%} \\
\trow \textbf{Unlevered rate $\rho$} & \textbf{11.38\%} \\
\arow Check: $\rho = \text{WACC} + K_d t (D/V)$ & ties\\
\end{tabular}
\end{extab}

\begin{extab}{Beta: bottom-up against regression}{Source: Damodaran January 2026;
60 monthly observations to March 2026 against the NIFTY 500.}
\begin{tabular}{@{}p{0.60\linewidth}r@{}}
\hrow \wh{Measure} & \wh{Value}\\
Bottom-up levered beta (used) & 0.908\\
\arow Regression beta, 60 monthly obs. & 1.151\\
R-squared & 0.267\\
\arow Blume-adjusted regression beta & 1.101\\
Difference, regression less bottom-up & $+$0.243\\
\end{tabular}
\end{extab}

\begin{notebox}\boxtitlea{Why the bottom-up beta is preferred}
The regression explains only 27\% of the variance and the price series was supplied
with a vendor adjusted-close column rather than in exchange bhavcopy format. A
bottom-up beta from 469 global building-materials firms is the more stable estimate.
Using the regression beta of 1.151 instead would raise the cost of equity to 13.13\%
and \emph{lower} the value to roughly Rs 71 — the direction is against the company,
so the conservative choice is the one we did not take.
\end{notebox}

\switchcolumn*
\section{Discounted cash flow valuation}

\subsection{The cash flow profile}

Free cash flow to the firm is negative in FY2026-27 and FY2027-28. That is not a
modelling error. It is what a company running at 67.9\% utilisation with Rs 9,085
crore of capital work in progress still to commission looks like: capital is going
out faster than the assets already built are bringing it in. FCFF turns positive in
FY2028-29 and reaches Rs 6,693 crore by FY2035-36.

\subsection{The terminal year}

Terminal growth is 6.99\%, built multiplicatively as 65\% of nominal GDP, where
nominal GDP is $(1+6.5\%)(1+4.0\%)-1 = 10.76\%$ — not 10.5\%. The same fraction is
applied to both companies in this document so the two are treated consistently.
Terminal reinvestment equals $g$ times \emph{operating} invested capital, excluding
goodwill, because organic growth does not require buying more goodwill; charging
$g$ on goodwill as well would demand that the company keep making acquisitions
forever merely to stand still.

\subsection{From enterprise value to the owners' claim}

The bridge deducts every claim that ranks ahead of, or alongside, the ordinary
shareholder. Two items deserve note. Deferred tax liabilities of Rs 3,466 crore are
deducted once at the valuation date rather than amortised: a deferred tax liability
means cash tax has lagged book tax, so as it unwinds free cash flow falls, and taking
it as a single deduction avoids pretending to know the timing. Non-controlling
interest of Rs 12,499 crore is deducted because consolidated cash flows include 100\%
of subsidiaries the parent does not wholly own; omitting it would overstate value by
about Rs 51 per share.

\begin{keybox}\boxtitle{Result}
Enterprise value Rs 39,833 crore. Equity value attributable to owners Rs 24,575
crore. Value per share \textbf{Rs 99.4} against a traded Rs 400.90, a downside of
75.2\%.
\end{keybox}

\switchcolumn
\ex{amb_fcff.png}{\textbf{Exhibit A3.} FCFF equals NOPAT plus depreciation less
capex less the increase in net working capital. Source: authors' model.}

\ex{amb_bridge.png}{\textbf{Exhibit A2.} Every liability included in the cost of
capital is subtracted here, and nothing is subtracted twice. Source: authors' model;
FY2025-26 consolidated balance sheet.}

\ex{amb_tv.png}{\textbf{Exhibit A14.} A ten-year explicit horizon holds terminal
value to 72\% of enterprise value, below the 80--90\% threshold at which the
terminal assumption effectively becomes the valuation. Source: authors' model.}

\begin{extab}{DCF summary}{Source: authors' model.}
\begin{tabular}{@{}p{0.58\linewidth}r@{}}
\hrow \wh{Line} & \wh{Rs cr}\\
PV of explicit FCFF, FY27--FY36 & 11,060\\
\arow PV of terminal value & 28,773\\
\trow \textbf{Enterprise value of operations} & \textbf{39,833}\\
Add: cash, bank and current investments & 959\\
\arow Less: borrowings and leases & $-$866\\
Add: tax paid under protest & 614\\
\arow Less: deferred tax liabilities & $-$3,466\\
Less: non-controlling interest & $-$12,499\\
\trow \textbf{Equity value, owners} & \textbf{24,575}\\
Shares outstanding (crore) & 247.18\\
\trow \textbf{Value per share (Rs)} & \textbf{99.4}\\
\arow Terminal value \% of EV & 72.3\% \\
Implied terminal EV/EBITDA & 4.65x\\
\end{tabular}
\end{extab}

\switchcolumn*
\section{Method equivalence and consistency audit}

\subsection{All five methods agree}

FCFF discounted at WACC, capital cash flow and adjusted present value discounted at
the unlevered rate $\rho$, FCFE discounted at the cost of equity, and economic value
added built from opening invested capital all return the same equity value per share.
The maximum difference across the five is zero to six decimal places.

\begin{okbox}\boxtitleg{What agreement does and does not prove}
Agreement across five methods proves internal consistency and nothing more. It says
nothing whatever about whether the assumptions are right. We were also able to
reproduce the entire model independently in Python from the input assumptions alone,
and both the value per share and the EVA-to-DCF identity tie exactly.
\end{okbox}

\subsection{The five consistencies}

\subsubsection{Assumptions --- \PASS} Volume is capacity times utilisation, so
growth is tied to physical assets. Capex is an explicit input anchored to the
expansion plan and is never held flat while volume rises. The terminal year is forced
to steady state with reinvestment equal to $g$ times operating invested capital,
which is the condition that makes EVA reconcile.

\subsubsection{Risk --- \PASS} Firm flows at firm rates, equity flows at the cost of
equity. Constant $D/V$ rebalancing is assumed, so the tax shield is discounted at
$\rho$, and the relevering formula carries no tax term to match. The identity $\rho =
\text{WACC} + K_d t (D/V)$ ties exactly.

\subsubsection{Asset --- \PART} Other income is excluded from EBITDA, so cash and
investments are non-operating and added back once. The unlevered beta is
cash-corrected for the same reason. \emph{The gap:} depreciation is charged at 8.0\%
of operating fixed assets, but the historical charge includes amortisation of the
Rs 9,433 crore intangible balance, which the model holds flat and never amortises.
The charge and the base do not cover the same assets. Splitting the two requires a
disclosed split that was not obtained.

\subsubsection{Liability --- \PASS} Borrowings and leases sit in the WACC weights and
in the bridge. Deferred tax is in neither the weights nor invested capital, because
it bears no interest, but is deducted in the bridge. Non-controlling interest is
deducted. Nothing appears twice and nothing is missing.

\subsubsection{Multiplier --- \PART} The implied terminal EV/EBITDA of 4.65x sits far
below the cement peer median of 19.4x. The two do not bracket. We take the position
that the peer multiple is the implausible end — it requires perpetual growth above
nominal GDP — but the gap is real and is recorded as a PARTIAL rather than smoothed
over.

\switchcolumn
\begin{extab}{Method equivalence}{All five methods are built from the same forecast
and must agree; any difference would be an error, not a modelling choice. Source:
authors' model.}
\begin{tabular}{@{}p{0.30\linewidth}p{0.22\linewidth}rr@{}}
\hrow \wh{Method} & \wh{Rate} & \wh{Rs/sh} & \wh{Diff}\\
FCFF & WACC & 99.42 & 0.00\\
\arow CCF & $\rho$ & 99.42 & 0.00\\
APV & $\rho$ & 99.42 & 0.00\\
\arow FCFE & $K_e$ & 99.42 & 0.00\\
EVA & WACC & 99.42 & 0.00\\
\trow \textbf{Status} & & \multicolumn{2}{r}{\textbf{RECONCILED}}\\
\end{tabular}
\end{extab}

\begin{extab}{Mohanty Appendix 5A: the seven sanity checks}{Source: authors'
assessment against the course audit framework.}
\begin{tabular}{@{}p{0.60\linewidth}l@{}}
\hrow \wh{Check} & \wh{Verdict}\\
1. Revenue growth plausible vs GDP & \PASS\\
\arow 2. Capacity supports projected volume & \PASS\\
3. Capex sufficient to build capacity & \PASS\\
\arow 4. Margin expansion realism & \PART\\
5. Working capital sanity & \PASS\\
\arow 6. Terminal ROIC fades toward WACC & \PASS\\
7. Terminal value dominance & \PASS\\
\end{tabular}
\end{extab}

\begin{notebox}\boxtitlea{Check 4, and the direction of the error}
EBITDA per tonne rises 70\% over ten years to Rs 1,502, which is 56\% above what
UltraTech's best assets earn today. That is demanding. It is demanding
\emph{against us}: a reader who regards Rs 1,502 as unreachable should note that the
value would be lower still, and the SELL stronger.
\end{notebox}

\begin{okbox}\boxtitleg{Check 6 passes here}
Terminal ROIC on operating capital is 11.94\% against a WACC of 11.36\%, a spread of
0.58 percentage points. That is a defensible steady state for a commodity producer
and assumes no perpetual supernormal rent. On invested capital \emph{including}
goodwill the terminal return remains below the cost of capital.
\end{okbox}

\switchcolumn*
\section{Relative valuation}

Multiples are computed after the discounted cash flow, never before it, and are used
to challenge the DCF rather than to set it. On the peer set Ambuja looks cheap: at
17.0x EV/EBITDA it is the lowest of the five names against a peer median of 19.4x,
and at 1.67x price-to-book it is the lowest by a wide margin.

\subsection{The two approaches disagree, and the disagreement is the finding}

Relative valuation makes Ambuja look roughly a fifth cheap. The DCF says it is worth
a quarter of its price. The resolution is in the price-to-book row. EV/EBITDA cannot
see the Rs 22,979 crore of goodwill Ambuja capitalised on its acquisitions; price to
book can. Across the peer set, price to book tracks return on equity closely — JK
Cement at 5.6x on a 16\% ROE, Shree at 3.7x on 7\%, Ambuja at 1.7x on 3.8\%. A
company earning a lower return \emph{should} trade at a lower multiple of book.

\begin{keybox}\boxtitle{Cheap is the diagnosis, not the opportunity}
Ambuja is priced at about 79\% of UltraTech's EV/EBITDA, so it looks a fifth cheaper.
It is priced at about 43\% of UltraTech's enterprise value per rupee of invested
capital. But it earns only about 31\% of UltraTech's return on that capital. The
discount is nowhere near large enough for the gap in returns.
\end{keybox}

\subsection{A caution on the price-to-book row}

Applying a peer price-to-book multiple to a company whose return on equity differs
sharply from the peer set is not a valuation; it is a category error. The implied
values of Rs 826--1,102 are shown for completeness, and because the pairing is itself
the diagnostic, but they should not be carried into a conclusion without adjusting
for the ROE difference.

\subsection{Basis warning}

Ambuja is carried on a consolidated basis and every peer on a standalone basis. This
is necessary rather than sloppy: Ambuja's market capitalisation prices the whole
group while its standalone EBITDA excludes ACC entirely. For the other four companies
standalone and consolidated differ by under 10\%, so the problem does not arise.
Ambuja is also one quarter behind the others, on FY2025-26 rather than trailing
twelve months to June 2026.

\switchcolumn
\ex{amb_peers.png}{\textbf{Exhibit A11.} Ambuja on a consolidated basis with
non-controlling interest added to enterprise value; peers standalone. Source: peer
annual reports and disclosed market capitalisations, trailing twelve months to June
2026.}

\ex{amb_pbroe.png}{\textbf{Exhibit A12.} The most important exhibit in this section.
Price to book tracks return on equity across the peer set; Ambuja sits at the bottom
of both axes. Source: peer annual reports; authors' analysis.}

\begin{extab}{Peer set, FY2025-26 / TTM June 2026}{Ambuja consolidated with NCI in
EV; peers standalone. Source: company annual reports and disclosed market
capitalisations.}
\begin{tabular}{@{}p{0.30\linewidth}rrrr@{}}
\hrow \wh{Company} & \wh{EV/EB} & \wh{P/E} & \wh{P/B} & \wh{ROE}\\
\textbf{Ambuja} & \textbf{17.0} & 43.9 & \textbf{1.67} & \textbf{3.8\%}\\
\arow UltraTech & 21.1 & 40.8 & 4.24 & 10.4\% \\
Shree Cement & 19.1 & 54.2 & 3.69 & 6.8\% \\
\arow JK Cement & 19.7 & 35.8 & 5.64 & 15.7\% \\
Ramco Cements & 19.0 & 122.9 & 2.69 & 2.2\% \\
\trow \textbf{Peer median} & \textbf{19.4} & \textbf{47.5} & \textbf{3.97} & \textbf{8.6\%}\\
\arow Ambuja, DCF-implied & 6.2 & & & \\
\end{tabular}
\end{extab}

\begin{warnbox}\boxtitler{The sector cannot be the benchmark}
Our own model says UltraTech — the largest member of this peer median — is 53\%
overvalued. A median built from such companies cannot be used to argue that anything
in the set is cheap. We are not claiming Ambuja is mispriced against its sector. We
are claiming the sector is mispriced against the cash it generates, and that Ambuja
is the member with the weakest returns.
\end{warnbox}

\switchcolumn*
\section{Sensitivity, scenarios and catalysts}

\subsection{Read the direction carefully}

The WACC-against-growth grid slopes \emph{upward} with terminal growth: at an 11\%
WACC the value rises from Rs 95 at 2\% growth to Rs 108 at 6\%. That is because this
model assumes return on capital climbs to 11.94\% by the terminal year, just above
the cost of capital, so terminal growth adds value. In the early years, when the
company earns roughly 7\% against a cost of 11.4\%, growth destroys value.

The real diagnostic is therefore the crossing point. Everything depends on whether
return on capital actually climbs through the cost of capital. That is an assumption,
not an observation, and Ambuja has not earned above its cost of capital in either
year of consolidated history available.

\subsection{What the traded price requires}

Solved backwards, Rs 400.90 requires a weighted average cost of capital of 8.60\%
against our 11.36\%. For a company funded 99\% by equity that implies an equity risk
premium far below any defensible figure. Holding the cost of capital and solving on
operations instead requires a terminal return on capital comfortably above 15\%,
EBITDA per tonne above Rs 2,000, and utilisation above 90\%.

\subsection{Catalysts}

A gap of this size cannot be presented as mispricing without naming what would close
it. \emph{On the bull side:} quarterly utilisation and EBITDA-per-tonne prints
confirming that the Sanghi, Penna, Orient and ACC integrations are lifting blended
unit economics; completion of the merger, which removes holding-company complexity
and the minority leakage; and a cement price recovery. \emph{On the bear side:}
industry capacity additions outrunning demand and holding utilisation near 70\%,
which would keep return on capital pinned below its cost indefinitely.

The single most informative disclosure to watch is consolidated EBITDA per tonne,
because it is the one number that separates this valuation from the market price.

\begin{notebox}\boxtitlea{On forecasting this company}
A sell-side house modelled FY2025-26 EBITDA of Rs 8,477 crore in December 2025, nine
months into the year. Ambuja reported Rs 6,539 crore — a 30\% miss — on a revenue
estimate accurate to within 2.7\%. The entire error was margin. Treat any forecast of
rapid margin recovery with that record in mind.
\end{notebox}

\switchcolumn
\ex{amb_sens.png}{\textbf{Exhibit A10.} Live two-way grid recomputed from the
forecast cash flows and the steady-state terminal formula. No combination of a
defensible WACC and a terminal growth rate below it reaches Rs 401. Source: authors'
model.}

\ex{amb_reverse.png}{\textbf{Exhibit A16.} Indexed so that the model base case equals
1.0 on each driver. Source: authors' model.}

\begin{extab}{Scenario view}{Scenarios are tied to driver assumptions, not to
arbitrary percentage shifts. Source: authors' model.}
\begin{tabular}{@{}p{0.34\linewidth}rrr@{}}
\hrow \wh{Driver} & \wh{Bear} & \wh{Base} & \wh{Bull}\\
Utilisation by FY36 & 75\% & 85\% & 90\% \\
\arow EBITDA/t by FY36 (Rs) & 1,150 & 1,502 & 1,850\\
Terminal ROIC & 8.5\% & 11.9\% & 15.0\% \\
\arow Terminal growth & 5.0\% & 6.99\% & 7.5\% \\
\trow \textbf{Value per share (Rs)} & \textbf{$\sim$45} & \textbf{99} & \textbf{$\sim$210}\\
\arow Gap to Rs 400.90 & $-$89\% & $-$75\% & $-$48\% \\
\end{tabular}
\end{extab}

\begin{warnbox}\boxtitler{Even the bull case does not get there}
On assumptions we would struggle to defend — 90\% utilisation, Rs 1,850 per tonne, a
15\% terminal return on capital — the value is roughly Rs 210, still 48\% below the
traded price.
\end{warnbox}

\switchcolumn*
\section{Risks and what would change our mind}

\subsection{Risks to the SELL}

\subsubsection{Integration delivers faster than modelled} If the Sanghi and Penna
assets reach parent-level margins in three years rather than ten, the standalone
margin path is too slow and the value rises materially.

\subsubsection{A genuine pricing cycle} Our realisation assumption grows below
inflation throughout. A sustained industry price recovery, which the sector has seen
before, would break that assumption in the company's favour.

\subsubsection{The merger unlocks more than the share count} If consolidation
delivers real cost synergies rather than only removing minority leakage, the
post-merger entity earns more than the sum we have valued.

\subsubsection{Our beta is too low} We use a bottom-up beta of 0.908 against a
regression beta of 1.151. We chose the lower figure, which \emph{raises} the value.
The risk runs the other way here.

\subsection{Limitations we accept}

Consolidated history is two years deep, so the return trend is asserted from two
points. The depreciation base and the depreciation charge do not cover the same
assets. The merger terms require re-sourcing from the scheme filing. And the implied
terminal multiple sits far below anything cement has traded at, which we have taken a
position on rather than resolved.

\begin{keybox}\boxtitle{What would change our mind}
Two consecutive years of consolidated ROIC including goodwill above 9\%, accompanied
by utilisation above 80\% and EBITDA per tonne above Rs 1,200. That combination would
demonstrate the crossing this valuation treats as an assumption, and would move the
value toward the upper half of our scenario range.
\end{keybox}

\switchcolumn
\begin{extab}{Recommendation summary}{Source: authors' model; NSE close 30 March
2026.}
\begin{tabular}{@{}p{0.58\linewidth}r@{}}
\hrow \wh{Item} & \wh{Value}\\
Recommendation & \textbf{SELL}\\
\arow Intrinsic value per share & Rs 99.4\\
Market price, 30 March 2026 & Rs 400.90\\
\arow Downside & $-$75.2\% \\
Market capitalisation & Rs 99,095 cr\\
\arow Model equity value & Rs 24,575 cr\\
WACC & 11.36\% \\
\arow Terminal growth & 6.99\% \\
Explicit forecast horizon & 10 years\\
\arow Terminal value \% of EV & 72.3\% \\
Implied terminal EV/EBITDA & 4.65x\\
\arow WACC implied by market price & 8.60\% \\
\end{tabular}
\end{extab}

\begin{okbox}\boxtitleg{Classification: the difficult company}
Ambuja is the difficult-to-value name in this pair, and the reasons are specific:
only two years of consolidated history, a restated prior year and a fifteen-month
transition year, a reported tax line that is a net credit making the effective rate
minus 71\%, a merger in progress that changes the share base, and a capital structure
in which a third of invested capital is goodwill. Each of those had to be handled
explicitly before any return measure became meaningful.
\end{okbox}
\end{paracol}
\clearpage
